Dominique répond à un QCM comportant 10 questions.

Pour chaque question, il est proposé 4 réponses dont une seule est exacte.

Dominique répond au hasard à chacune des 10 questions en cochant, pour chaque question, exactement une case parmi les 4.

Pour chacune des questions, la probabilité qu'il réponde correctement est donc $\dfrac{1}{4}$.

On note $X$ la variable aléatoire qui compte le nombre de bonnes réponses à ce QCM.

\begin{enumerate}
\item Déterminer la loi suivie par la variable aléatoire $X$ et donner les paramètres de cette loi.

\item Quelle est la probabilité que Dominique obtienne exactement 5 bonnes réponses ? Arrondir le résultat à $10^{-4}$ près.

\item Donner l'espérance de $X$ et interpréter ce résultat dans le contexte de l'exercice.

\item On suppose dans cette question qu'une bonne réponse rapporte un point et qu'une mauvaise réponse fait perdre 0,5 point. La note finale peut donc être négative.

On note $Y$ la variable aléatoire qui donne le nombre de points obtenus.

	\begin{enumerate}
		\item Calculer $P(Y=10)$, on donnera la valeur exacte du résultat.

		\item À partir de combien de bonnes réponses la note finale de Dominique est-elle positive ? Justifier.

		\item Calculer $P(Y \leqslant 0)$, on donnera une valeur approchée au centième.

		\item Montrer que $Y=1,5 X-5$.

		\item Calculer l'espérance de la variable aléatoire $Y$.
	\end{enumerate}
\end{enumerate}
